%!TEX root = ../main.tex
% =============================================================
%  CAPÍTULO 5 — JUSTIFICACIÓN ECONÓMICA (Entrega 5 · 15 % · 15/9)
% =============================================================

\chapter{Justificación económica}\label{cap:economica}

\begin{guia}[Guía de la cátedra: qué se evalúa en este capítulo]
Es el capítulo de mayor peso individual (15\,\%). Se evalúa como una
evaluación de proyectos de inversión:
\begin{enumerate}
    \item \textbf{Análisis de mercado y competidores}: tamaño, segmentos,
          precios de referencia.
    \item \textbf{Modelo de negocio}: B2B, B2C, B2G o interno; niveles de
          precio y fuentes de ingreso.
    \item \textbf{Inversión inicial} desglosada por rol y tiempo, más
          infraestructura.
    \item \textbf{Costos operativos} recurrentes.
    \item \textbf{Escenarios} pesimista, base y optimista, con flujo de caja
          proyectado a 3 a 5 años.
    \item \textbf{Indicadores}: \gls{van}, \gls{tir} y período de repago,
          con una tasa de descuento construida y justificada con fuentes.
    \item \textbf{Supuestos y restricciones} explícitos.
\end{enumerate}
Todo número con su origen: cotización, tarifa pública, estadística
citada o estimación propia declarada como tal.
\end{guia}

\section{Introducción y objetivo del capítulo}

El objetivo de este capítulo es presentar los beneficios tangibles e
intangibles, desde el punto de vista económico, que se esperan lograr con
el desarrollo e implementación de la plataforma inteligente de gestión
integral de eventos propuesta en los capítulos anteriores.

Para esto se realiza un análisis costo-beneficio, se calculan los
principales indicadores de rentabilidad (flujo de fondos, período de
repago, \gls{van} y \gls{tir}) y se presenta un análisis de impacto
institucional, legal, ambiental y social, junto con el caso de negocio
que justifica la alternativa seleccionada frente a otras posibles.

Todos los montos de este capítulo se expresan en dólares estadounidenses
(USD), utilizando como referencia la cotización del dólar oficial de
Argentina de septiembre de 2026 (USD 1 = ARS 1.500), con el fin de evitar
que la alta inflación local distorsione la comparación de los valores a
lo largo de los cinco años proyectados.

\section{Estimación de la cantidad de usuarios y contrataciones}

Dado que las encuestas y entrevistas de la Fase 2 de la metodología
todavía se encuentran en curso, la
cantidad de usuarios y de contrataciones utilizada en este capítulo se
estimó a partir de fuentes secundarias sobre el mercado de eventos en
Argentina, como una primera aproximación que deberá ajustarse una vez
procesados los resultados reales de las encuestas.

Como referencia de mercado se relevó que:

\begin{itemize}
    \item \textbf{Alcance de mercado}: los principales directorios y
          marketplaces de salones que operan en Argentina (InEventos,
          Casamientos.com.ar, entre otros) listan más de 100 salones a
          nivel nacional, con una concentración importante en CABA y el
          Gran Buenos Aires (más de 40 salones solo en CABA).
    \item \textbf{Volumen por salón}: según la Asociación de
          Organizadores de Fiestas, Reuniones y Eventos Empresariales y
          sus Proveedores de la República Argentina (Aoefrep), un salón
          promedio realiza alrededor de 80 eventos por año.
\end{itemize}

A partir de estos datos se optó por una estimación deliberadamente
conservadora para el lanzamiento: en el Año 1 la plataforma no busca
capturar el volumen completo de un salón promedio, sino operar como una
prueba piloto acotada, del orden de 20 contrataciones gestionadas a
través del marketplace, creciendo de forma gradual a medida que se suman
salones y proveedores. La plataforma se lanza inicialmente en Buenos
Aires y el Gran Buenos Aires, la misma zona donde opera Rodríguez
Catering, el caso real utilizado en el prototipo.

\begin{table}[H]
    \centering
    \caption{Evolución estimada de la adopción}
    \label{tab:adopcion}
    \small
    \begin{tabularx}{\textwidth}{L R R R R R}
        \toprule
        \tb{Concepto} & \tb{Año 1} & \tb{Año 2} & \tb{Año 3} & \tb{Año 4} & \tb{Año 5} \\
        \midrule
        Salones/proveedores suscriptos (acum.) & 15 & 40 & 90 & 160 & 250 \\
        Contrataciones vía marketplace & 20 & 75 & 200 & 400 & 750 \\
        \bottomrule
    \end{tabularx}
\end{table}

Esta proyección implica que recién hacia el Año 5 la plataforma
alcanzaría una penetración de entre el 20\,\% y el 30\,\% del total de
salones identificados en los directorios relevados a nivel nacional, lo
cual se considera una meta ambiciosa pero razonable para un proyecto de
este tipo dentro del horizonte analizado.

\section{Análisis costo-beneficio}

\subsection{Costos de inversión inicial (Año 0)}

Corresponden a los costos necesarios para llevar la plataforma, ya
construida sobre el caso real de Rodríguez Catering, a una primera
versión publicada y operativa.

\begin{table}[H]
    \centering
    \caption{Inversión inicial desglosada}
    \label{tab:inversion}
    \small
    \begin{tabularx}{\textwidth}{c L r}
        \toprule
        \tb{\#} & \tb{Descripción del costo} & \tb{Monto (USD)} \\
        \midrule
        1 & Desarrollo de software & 6.000 \\
        2 & Diseño de interfaz (\gls{ux}/UI) y branding de la plataforma & 530 \\
        3 & Infraestructura inicial (servidor cloud, dominio, certificados SSL) & 200 \\
        4 & Licencias e integraciones de terceros (pagos, mapas, mensajería) & 170 \\
        5 & Marketing de lanzamiento & 430 \\
        \midrule
        \multicolumn{2}{r}{\tb{Total inversión inicial}} & \tb{7.330} \\
        \bottomrule
    \end{tabularx}
\end{table}

El monto de inversión resulta comparativamente bajo porque el desarrollo
inicial es realizado por el propio tesista como parte del proyecto, y
porque los costos de infraestructura y desarrollo en Argentina son
considerablemente menores en dólares que en otros países de la región,
dado el tipo de cambio actual.

\subsection{Costos operativos (post-lanzamiento)}

\begin{table}[H]
    \centering
    \caption{Costos operativos proyectados, en USD}
    \label{tab:costos-operativos}
    \small
    \begin{tabularx}{\textwidth}{L R R R R R}
        \toprule
        \tb{Concepto} & \tb{Año 1} & \tb{Año 2} & \tb{Año 3} & \tb{Año 4} & \tb{Año 5} \\
        \midrule
        Hosting y servicios cloud & 400 & 600 & 930 & 1.470 & 2.130 \\
        Mantenimiento y soporte técnico & 1.200 & 1.470 & 1.870 & 2.400 & 3.200 \\
        Marketing y adquisición de usuarios & 800 & 1.200 & 1.800 & 2.800 & 4.000 \\
        Atención al cliente / community mgmt. & 600 & 730 & 1.400 & 2.630 & 4.000 \\
        \midrule
        \tb{Total costos operativos} & \tb{3.000} & \tb{4.000} & \tb{6.000} & \tb{9.300} & \tb{13.300} \\
        \bottomrule
    \end{tabularx}
\end{table}

\subsection{Beneficios (fuentes de ingreso)}

De acuerdo con el modelo de negocio definido para la plataforma, se
adopta un modelo mixto de ingresos, compuesto por:

\begin{itemize}
    \item \textbf{Suscripción mensual base}: para que salones y
          proveedores puedan administrar reservas, presupuestos,
          contratos y pagos, y aparecer en el marketplace. Para
          incentivar la adopción temprana, se fija un precio de
          lanzamiento reducido de USD 5/mes durante los Años 1 y 2, que
          pasa a USD 15/mes a partir del Año 3, una vez consolidada la
          base de usuarios.
    \item \textbf{Complemento de \gls{ia}}: opcional, con un costo
          adicional muy bajo de USD 5/mes, que habilita el acceso a las
          recomendaciones inteligentes y a las simulaciones del Gemelo
          Digital. Se mantiene deliberadamente accesible para
          incentivar su adopción y generar el volumen de datos de uso
          necesario para mejorar el motor de recomendación.
    \item \textbf{Comisión por contratación}: del 5\,\% sobre el valor
          del servicio contratado a través del marketplace, tomando un
          ticket promedio estimado de USD 1.500 por contratación.
\end{itemize}

\begin{table}[H]
    \centering
    \caption{Composición de ingresos proyectados}
    \label{tab:ingresos}
    \small
    \begin{tabularx}{\textwidth}{L R R R R R}
        \toprule
        \tb{Concepto} & \tb{Año 1} & \tb{Año 2} & \tb{Año 3} & \tb{Año 4} & \tb{Año 5} \\
        \midrule
        Salones/proveedores suscriptos & 15 & 40 & 90 & 160 & 250 \\
        Precio suscripción base/mes (USD) & 5 & 5 & 15 & 15 & 15 \\
        \% que contrata el complemento de IA & 20\,\% & 30\,\% & 45\,\% & 60\,\% & 70\,\% \\
        Ingreso promedio por suscriptor/mes (USD) & 6,0 & 6,5 & 17,3 & 18,0 & 18,5 \\
        Ingresos por suscripción (anual) & 1.080 & 3.120 & 18.630 & 34.560 & 55.500 \\
        \midrule
        Contrataciones vía marketplace & 20 & 75 & 200 & 400 & 750 \\
        Ticket promedio por contratación (USD) & 1.500 & 1.500 & 1.500 & 1.500 & 1.500 \\
        Ingresos por comisión (5\,\%) & 1.500 & 5.625 & 15.000 & 30.000 & 56.250 \\
        \midrule
        \tb{Total ingresos} & \tb{2.580} & \tb{8.745} & \tb{33.630} & \tb{64.560} & \tb{111.750} \\
        \bottomrule
    \end{tabularx}
\end{table}

\section{Flujo de fondos proyectado}

\begin{table}[H]
    \centering
    \caption{Flujo de fondos proyectado, en USD}
    \label{tab:flujo-caja}
    \small
    \begin{tabularx}{\textwidth}{l R R R R R R}
        \toprule
        \tb{} & \tb{Año 0} & \tb{Año 1} & \tb{Año 2} & \tb{Año 3} & \tb{Año 4} & \tb{Año 5} \\
        \midrule
        Inversión inicial & $-$7.330 & & & & & \\
        Ingresos & & 2.580 & 8.745 & 33.630 & 64.560 & 111.750 \\
        Costos operativos & & $-$3.000 & $-$4.000 & $-$6.000 & $-$9.300 & $-$13.300 \\
        \midrule
        \tb{Flujo de fondos neto} & \tb{$-$7.330} & \tb{$-$420} & \tb{4.745} & \tb{27.630} & \tb{55.260} & \tb{98.450} \\
        \tb{Flujo acumulado} & \tb{$-$7.330} & \tb{$-$7.750} & \tb{$-$3.005} & \tb{24.625} & \tb{79.885} & \tb{178.335} \\
        \bottomrule
    \end{tabularx}
\end{table}

\section{Indicadores de rentabilidad}

\subsection{Período de repago (\emph{payback})}

Con el precio de suscripción reducido de los primeros dos años, el
flujo neto del Año 1 pasa a ser levemente negativo, y el flujo
acumulado recién se vuelve positivo durante el Año 3. Interpolando
dentro de ese año:
\[
    \text{Payback} \approx 2 \text{ años} + \frac{3.005}{27.630}
    \approx 2{,}1 \text{ años (aproximadamente 2 años y 1 mes)}.
\]

\subsection{Valor actual neto (\acrshort{van})}

Aplicando una tasa de corte del 20\,\% anual (se utiliza una tasa
relativamente alta para reflejar tanto el riesgo propio de un proyecto
tecnológico en etapa temprana como el riesgo país de Argentina):
\begin{equation}
    \mathit{VAN} = \sum_{t=0}^{5} \frac{\gls{sym:fc}}{(1+\gls{sym:r})^{t}}
    \approx \text{USD } 77.800,
    \label{eq:van}
\end{equation}
donde $r = 0{,}20$. Al ser $\gls{van} \geq 0$, según el criterio de
decisión, el proyecto es aceptado.

\subsection{Tasa interna de retorno (\acrshort{tir})}

Calculando la tasa que hace al \gls{van} igual a cero sobre el mismo
flujo de fondos, se obtiene una \gls{tir} aproximada del 122\,\%.

Este valor se explica principalmente por el bajo monto de inversión
inicial (USD 7.330), consecuencia de que gran parte del desarrollo se
realiza como trabajo propio del tesista y no como una contratación de
mercado. Un análisis más conservador podría incorporar un costo de
oportunidad por el tiempo dedicado por el desarrollador, lo cual
reduciría la \gls{tir} sin modificar la conclusión general: el proyecto
resulta económicamente viable bajo los supuestos planteados.

\subsection{Escenarios pesimista y optimista}

Todo lo desarrollado hasta acá corresponde al escenario base. Para
completar el análisis se construyen dos escenarios alternativos,
aplicando supuestos distintos de adopción, de velocidad del aumento de
precios y de costos sobre la misma estructura del escenario base.

\begin{table}[H]
    \centering
    \caption{Supuestos por escenario}
    \label{tab:supuestos-escenarios}
    \small
    \begin{tabularx}{\textwidth}{L L L L}
        \toprule
        \tb{Supuesto} & \tb{Pesimista} & \tb{Base} & \tb{Optimista} \\
        \midrule
        Adopción (salones/proveedores y contrataciones) & 60\,\% de la proyección base & Tabla~\ref{tab:adopcion} & 150\,\% de la proyección base \\
        Aumento del precio de suscripción a USD 15/mes & Recién en el Año 4 & En el Año 3 & Ya en el Año 2 \\
        Adopción del complemento de \gls{ia} & 15\,\% a 50\,\% (Años 1--5) & 20\,\% a 70\,\% (Años 1--5) & 25\,\% a 80\,\% (Años 1--5) \\
        Costos operativos & 85\,\% de los costos base & Tabla~\ref{tab:costos-operativos} & 120\,\% de los costos base (más soporte y marketing) \\
        \bottomrule
    \end{tabularx}
\end{table}

\begin{table}[H]
    \centering
    \caption{Flujo de fondos neto por escenario, en USD}
    \label{tab:flujo-escenarios}
    \small
    \begin{tabularx}{\textwidth}{l R R R R R R}
        \toprule
        \tb{Escenario} & \tb{Año 0} & \tb{Año 1} & \tb{Año 2} & \tb{Año 3} & \tb{Año 4} & \tb{Año 5} \\
        \midrule
        Pesimista & $-$7.330 & $-$1.029 & 1.703 & 8.112 & 29.679 & 53.945 \\
        Base      & $-$7.330 & $-$420   & 4.745 & 27.630 & 55.260 & 98.450 \\
        Optimista & $-$7.330 & 300      & 15.840 & 44.055 & 87.120 & 153.915 \\
        \bottomrule
    \end{tabularx}
\end{table}

En el escenario pesimista, el flujo acumulado recién se vuelve positivo
durante el Año 3 (más tarde que en el escenario base), pero el proyecto
sigue siendo viable: nunca se necesita una segunda ronda de inversión, y
el flujo acumulado a cinco años es positivo (USD 85.080). En el escenario
optimista, el flujo ya es positivo desde el Año 1.

\begin{table}[H]
    \centering
    \caption{Indicadores financieros por escenario}
    \label{tab:indicadores-financieros}
    \small
    \begin{tabular}{l r r r}
        \toprule
        \tb{Escenario} & \tb{VAN (USD)} & \tb{TIR} & \tb{Repago} \\
        \midrule
        Pesimista & 33.700 & 78\,\% & $\approx$2,8 años \\
        Base      & 77.800 & 122\,\% & $\approx$2,1 años \\
        Optimista & 133.300 & 169\,\% & $\approx$1,4 años \\
        \bottomrule
    \end{tabular}
\end{table}

Incluso en el escenario pesimista --- con seis de cada diez de los
salones/proveedores y contrataciones previstos, precios de lanzamiento
más prolongados y menor adopción del complemento de \gls{ia}--- el
\gls{van} sigue siendo positivo y la \gls{tir} supera ampliamente la
tasa de corte del 20\,\%. Esto sugiere que la conclusión de viabilidad
no depende de que se cumplan exactamente los supuestos más optimistas de
adopción.

\section{Análisis de impacto}

\subsection{Justificación institucional}

En la actualidad, tanto Rodríguez Catering como los demás salones y
proveedores relevados administran su información de forma manual o
mediante herramientas independientes (WhatsApp, Excel, planillas en
papel), lo que genera duplicación de tareas y pérdida de información. La
situación deseada es contar con una plataforma única que centralice la
gestión comercial y operativa, reduciendo el tiempo administrativo
dedicado a tareas repetitivas.

\subsection{Justificación legal}

La plataforma debe cumplir con la normativa vigente de protección de
datos personales (Ley 25.326 de Protección de Datos Personales) dado que
administra datos de clientes, proveedores y pagos, y con las
regulaciones aplicables a plataformas de intermediación comercial y
procesamiento de pagos electrónicos.

\subsection{Justificación ambiental}

El impacto ambiental directo es reducido, dado que se trata de un
desarrollo de software. Como impacto positivo indirecto, la
digitalización de contratos, presupuestos y comunicaciones reduce el uso
de papel respecto de los procesos actuales, muchos de los cuales aún se
gestionan mediante documentación impresa.

\subsection{Justificación social}

La plataforma busca facilitar el acceso a servicios de eventos tanto a
clientes como a pequeños proveedores y salones que actualmente dependen
de la difusión boca en boca o de redes sociales para conseguir clientes,
ampliando su visibilidad comercial sin necesidad de una gran inversión
propia en marketing digital. Además, el precio de suscripción reducido
durante los primeros dos años busca facilitar específicamente la
adopción por parte de proveedores más pequeños, que suelen tener menor
capacidad de inversión inicial en herramientas digitales.

\section{Caso de negocio: alternativas consideradas}

\begin{itemize}
    \item \textbf{Mantener el statu quo}: continuar utilizando
          herramientas independientes (redes sociales, WhatsApp,
          planillas de cálculo). Se descarta porque no resuelve la
          fragmentación de la información ni ofrece herramientas
          inteligentes de recomendación o simulación.
    \item \textbf{Alternativa A -- Marketplace simple sin herramientas de
          gestión}: desarrollar únicamente un buscador de proveedores
          similar a BIG Fiestas o FEX Eventos. Se descarta porque no
          resuelve la gestión posterior del evento ni diferencia al
          proyecto de las plataformas ya existentes en el mercado
          argentino.
    \item \textbf{Alternativa B -- Sistema de gestión (SaaS) sin
          marketplace}: desarrollar solamente una herramienta de
          administración para salones y proveedores, sin conectar con
          clientes. Se descarta porque no centraliza la contratación de
          servicios.
    \item \textbf{Alternativa seleccionada}: plataforma integral que
          combina marketplace, gestión comercial/operativa y
          funcionalidades inteligentes (recomendación y Gemelo Digital),
          por ser la única que atiende de manera conjunta las
          necesidades identificadas en los tres grupos de usuarios
          (clientes, salones y proveedores).
\end{itemize}

\section{Conclusión del capítulo}

El análisis económico realizado, aun con carácter estimativo, muestra que
el proyecto resulta viable desde el punto de vista financiero en los tres
escenarios considerados: incluso en el escenario pesimista el \gls{van}
es positivo (USD 33.700) y la \gls{tir} (78\,\%) supera la tasa de corte.
En el escenario base, el más probable, se recupera la inversión inicial
en aproximadamente 2 años, con un \gls{van} de USD 77.800 y una \gls{tir}
del 122\,\%. El esquema de precios de lanzamiento reducido
retrasa levemente el punto de equilibrio respecto de un escenario de
precios más agresivo, pero resulta coherente con una estrategia de
adopción gradual orientada a consolidar primero una base sólida de
usuarios antes de aumentar el precio. A esto se suman beneficios
intangibles relevantes, como la reducción de la fragmentación de
información en la industria de eventos y la mayor visibilidad comercial
para proveedores de menor escala.

Se reitera que la cantidad de usuarios y contrataciones utilizada en
este capítulo es una estimación preliminar basada en fuentes secundarias
sobre el mercado de eventos en Argentina, y debe ser reemplazada por los
datos reales relevados en las encuestas de la Fase 2 en cuanto estén
disponibles.
